\kchapter{Security}{role}

\begin{objectives}
\begin{itemize}
  \item the layers of Kubernetes security (the ``4 Cs'');
  \item authentication, authorisation with RBAC, and ServiceAccounts;
  \item hardening Pods with \kc{securityContext} and Pod Security Standards;
  \item admission policies with ValidatingAdmissionPolicy, Kyverno, and Gatekeeper;
  \item supply-chain security: scanning, signing, and minimal images;
  \item a practical hardening checklist.
\end{itemize}
\end{objectives}

\section{Defence in depth: the 4 Cs}

\begin{center}
\begin{tikzpicture}
  \foreach \rad/\col in {3.3/k8snavy, 2.5/k8sdeep, 1.7/k8sblue, 0.9/k8smid}
    \fill[\col] (0,0) circle (\rad);
  \node[text=white, font=\bfseries] at (0,2.85) {\faCloud\ Cloud / datacentre};
  \node[text=white, font=\bfseries] at (0,2.05) {\kicon{k8s-logo-white}\ Cluster};
  \node[text=white, font=\bfseries] at (0,1.25) {\faCube\ Container};
  \node[text=k8snavy, font=\bfseries] at (0,0) {\faCode\ Code};
\end{tikzpicture}
\end{center}

Each layer protects the one inside it. A perfectly configured cluster cannot
save vulnerable code, and secure code cannot survive an exposed cloud account.
This chapter focuses on the \emph{cluster} and \emph{container} layers.

\section{Who are you? Authentication}

Kubernetes has no ``user'' objects. Humans are authenticated externally:

\begin{itemize}
  \item \textbf{client certificates} (common in kubeadm/kind admin configs);
  \item \textbf{OIDC} tokens from an identity provider (Google, Entra ID, Okta, Keycloak);
  \item \textbf{cloud IAM} integrations (EKS access entries, GKE IAM, AKS Entra ID).
\end{itemize}

Workloads authenticate as \kicon{sa}~\kterm{ServiceAccounts} --- real
namespaced objects. Every Pod runs as one (\kc{default} if you do not choose),
and receives a short-lived, automatically rotated token.

\section{What may you do? RBAC}

\kterm{Role-Based Access Control} answers ``may \emph{this subject} perform
\emph{this verb} on \emph{this resource}?'' with four objects:

\begin{center}
\resizebox{\linewidth}{!}{%
\begin{tikzpicture}[
  o/.style={draw=k8sblue, fill=white, line width=1pt, rounded corners=5pt,
            minimum width=3.1cm, minimum height=1.9cm, align=center, font=\small\bfseries,
            text=k8snavy},
  >={Stealth[length=2.6mm]}]
  \node[o] (u)  at (0,0)    {\kimg[0.8cm]{user}\\User / Group\\[-2pt]\scriptsize\mdseries or ServiceAccount};
  \node[o] (rb) at (4.5,0)  {\kimg[0.8cm]{rb}\\RoleBinding\\[-2pt]\scriptsize\mdseries (namespace)};
  \node[o] (r)  at (9,0)    {\kimg[0.8cm]{role}\\Role\\[-2pt]\scriptsize\mdseries verbs on resources};
  \node[o] (crb) at (4.5,-2.6) {\kimg[0.8cm]{crb}\\ClusterRoleBinding\\[-2pt]\scriptsize\mdseries (whole cluster)};
  \node[o] (cr) at (9,-2.6) {\kimg[0.8cm]{c-role}\\ClusterRole\\[-2pt]\scriptsize\mdseries cluster-wide rules};
  \draw[->, thick, k8sblue] (rb) -- node[above, font=\scriptsize]{subjects} (u);
  \draw[->, thick, k8sblue] (rb) -- node[above, font=\scriptsize]{roleRef} (r);
  \draw[->, thick, k8sdeep] (crb.west) -| (u.south);
  \draw[->, thick, k8sdeep] (crb) -- (cr);
  \draw[->, thick, k8sorange, dashed] (rb.south east) -- node[right, font=\scriptsize, pos=0.4]{may also reference} (cr.north west);
\end{tikzpicture}}
\end{center}

\begin{yamlcode}{rbac-ci-deployer.yaml}
apiVersion: v1
kind: ServiceAccount
metadata:
  name: ci-deployer
  namespace: shop
---
apiVersion: rbac.authorization.k8s.io/v1
kind: Role
metadata:
  name: deployer
  namespace: shop
rules:
  - apiGroups: ["apps"]
    resources: ["deployments"]
    verbs: ["get", "list", "watch", "update", "patch"]
  - apiGroups: [""]                     # "" = the core API group
    resources: ["pods", "pods/log"]
    verbs: ["get", "list", "watch"]
---
apiVersion: rbac.authorization.k8s.io/v1
kind: RoleBinding
metadata:
  name: ci-deployer-binding
  namespace: shop
subjects:
  - kind: ServiceAccount
    name: ci-deployer
    namespace: shop
roleRef:
  kind: Role
  name: deployer
  apiGroup: rbac.authorization.k8s.io
\end{yamlcode}

\begin{bashcode}{terminal}
kubectl auth can-i update deployments -n shop --as=system:serviceaccount:shop:ci-deployer   # yes
kubectl auth can-i delete secrets -n shop --as=system:serviceaccount:shop:ci-deployer       # no
kubectl auth can-i --list --as=system:serviceaccount:shop:ci-deployer -n shop
\end{bashcode}

\begin{warningbox}
Be stingy. Avoid \kc{cluster-admin} bindings, wildcard \kc{"*"} verbs or
resources, and granting \kc{get/list} on Secrets (listing Secrets reveals their
contents). Remember that the ability to \emph{create Pods} in a namespace
effectively grants access to every Secret and ServiceAccount in it.
\end{warningbox}

Built-in ClusterRoles cover most needs: \kc{view}, \kc{edit}, \kc{admin}, and
\kc{cluster-admin}. Bind \kc{view} or \kc{edit} per namespace for developers.

\begin{protip}
Set \kc{automountServiceAccountToken: false} on Pods that never call the
Kubernetes API. Fewer tokens lying around means less to steal.
\end{protip}

\section{Hardening Pods: securityContext}

\begin{yamlcode}{hardened-pod.yaml}
spec:
  automountServiceAccountToken: false
  securityContext:                       # Pod level
    runAsNonRoot: true
    runAsUser: 10001
    runAsGroup: 10001
    fsGroup: 10001
    seccompProfile:
      type: RuntimeDefault               # block dangerous syscalls
  containers:
    - name: api
      image: ghcr.io/example/api@sha256:3b1f...   # pin by digest
      securityContext:                   # container level
        allowPrivilegeEscalation: false
        readOnlyRootFilesystem: true
        privileged: false
        capabilities:
          drop: ["ALL"]
      volumeMounts:
        - name: tmp
          mountPath: /tmp                # writable scratch, since root fs is read-only
  volumes:
    - name: tmp
      emptyDir: {}
\end{yamlcode}

\section{Pod Security Standards}

Kubernetes defines three standard profiles, enforced per namespace by the
built-in \textbf{Pod Security Admission} controller (it replaced the removed
PodSecurityPolicy):

\begin{center}
\ktablestart
\begin{tabularx}{\linewidth}{P{2.4cm} L}
\kheader \kth{Profile} & \kth{Effect} \\
\kc{privileged} & No restrictions. Only for trusted system components. \\
\kc{baseline} & Blocks known privilege escalations: privileged containers, host namespaces, hostPath, etc. \\
\kc{restricted} & Baseline plus best practices: non-root, drop all capabilities, seccomp, no privilege escalation. \\
\end{tabularx}
\end{center}

\begin{yamlcode}{restricted-namespace.yaml}
apiVersion: v1
kind: Namespace
metadata:
  name: shop
  labels:
    pod-security.kubernetes.io/enforce: restricted   # reject violating Pods
    pod-security.kubernetes.io/warn: restricted      # warn on kubectl apply
    pod-security.kubernetes.io/audit: restricted     # record in audit log
\end{yamlcode}

\section{Policy as code}

For rules beyond Pod security --- ``images must come from our registry'',
``every Deployment needs an \kc{owner} label'', ``no \kc{latest} tags'' --- use
admission policies:

\begin{itemize}
  \item \textbf{ValidatingAdmissionPolicy} --- built in (GA since 1.30), rules
        written in CEL, no extra software.
  \item \textbf{Kyverno} --- policies written as Kubernetes YAML; can validate,
        mutate, generate, and verify image signatures.
  \item \textbf{OPA Gatekeeper} --- policies in the Rego language.
\end{itemize}

\begin{yamlcode}{no-latest-tag.yaml}
apiVersion: admissionregistration.k8s.io/v1
kind: ValidatingAdmissionPolicy
metadata:
  name: no-latest-tag
spec:
  failurePolicy: Fail
  matchConstraints:
    resourceRules:
      - apiGroups: ["apps"]
        apiVersions: ["v1"]
        operations: ["CREATE", "UPDATE"]
        resources: ["deployments"]
  validations:
    - expression: >-
        object.spec.template.spec.containers.all(c,
          !c.image.endsWith(':latest') && c.image.contains(':'))
      message: "Images must use an explicit, non-latest tag."
---
apiVersion: admissionregistration.k8s.io/v1
kind: ValidatingAdmissionPolicyBinding
metadata:
  name: no-latest-tag
spec:
  policyName: no-latest-tag
  validationActions: ["Deny"]
\end{yamlcode}

\section{Supply-chain security}

\begin{itemize}
  \item \textbf{Scan images} for known vulnerabilities in CI (Trivy, Grype)
        and continuously in the registry.
  \item \textbf{Use minimal base images} (distroless, Chainguard, \kc{-slim}).
  \item \textbf{Sign images} with Sigstore \textbf{cosign} and verify signatures
        at admission (Kyverno, Sigstore policy-controller).
  \item \textbf{Generate SBOMs} (software bills of materials) with Syft or BuildKit.
  \item \textbf{Pin by digest} so the image you tested is the image you run.
\end{itemize}

\begin{bashcode}{terminal}
trivy image --severity HIGH,CRITICAL ghcr.io/example/api:1.4.0
cosign sign ghcr.io/example/api@sha256:3b1f...
cosign verify ghcr.io/example/api@sha256:3b1f... \
  --certificate-identity-regexp 'https://github.com/example/.*' \
  --certificate-oidc-issuer https://token.actions.githubusercontent.com
\end{bashcode}

\section{Runtime security and auditing}

\begin{itemize}
  \item \textbf{Falco} or \textbf{Tetragon} detect suspicious behaviour at
        runtime (a shell spawned in a container, unexpected file writes).
  \item \textbf{Audit logging} on the API server records who did what and when.
  \item \textbf{kube-bench} checks a cluster against the CIS Kubernetes Benchmark.
\end{itemize}

\section{Hardening checklist}

\begin{multicols}{2}
\begin{itemize}
  \item[\textcolor{k8sgreen}{\faCheckSquare}] API server not public, or behind IP allow-lists
  \item[\textcolor{k8sgreen}{\faCheckSquare}] SSO/OIDC for humans; no shared admin kubeconfigs
  \item[\textcolor{k8sgreen}{\faCheckSquare}] Least-privilege RBAC, reviewed regularly
  \item[\textcolor{k8sgreen}{\faCheckSquare}] \kc{restricted} Pod Security on app namespaces
  \item[\textcolor{k8sgreen}{\faCheckSquare}] Default-deny NetworkPolicies
  \item[\textcolor{k8sgreen}{\faCheckSquare}] Secrets encrypted at rest; external secret manager
  \item[\textcolor{k8sgreen}{\faCheckSquare}] Images scanned, signed, pinned by digest
  \item[\textcolor{k8sgreen}{\faCheckSquare}] Non-root, read-only root FS, all capabilities dropped
  \item[\textcolor{k8sgreen}{\faCheckSquare}] Resource requests/limits everywhere
  \item[\textcolor{k8sgreen}{\faCheckSquare}] Audit logs shipped off-cluster
  \item[\textcolor{k8sgreen}{\faCheckSquare}] Cluster and nodes kept on supported versions
  \item[\textcolor{k8sgreen}{\faCheckSquare}] etcd backups tested by restoring
\end{itemize}
\end{multicols}

\begin{recap}
\begin{itemize}
  \item Authentication is external; ServiceAccounts identify workloads.
  \item RBAC = Roles (what) + Bindings (who); grant the minimum.
  \item Harden Pods with \kc{securityContext}; enforce with Pod Security Admission.
  \item Use admission policies, image scanning, and signing to secure the supply chain.
\end{itemize}
\end{recap}
